\problem{P18}{随机排列的短循环}{组合计数与阶乘矩}
在 $n$ 个元素的均匀随机排列中，记 $C_k$ 为长度恰为 $k$ 的循环数，$1\le k\le n$。求 $\E C_k$、$\Var C_k$，以及 $j\ne k$ 时的 $\Cov(C_j,C_k)$。说明长度受到总元素数限制时，为什么协方差会改变。
\solution 对每个 $k$ 元子集 $A$，令 $I_A$ 表示这些元素恰好构成一个循环。其内部循环有 $(k-1)!$ 种，剩余元素任意排列，所以
\[
\Prob(I_A=1)=\frac{(k-1)!(n-k)!}{n!},
\qquad \boxed{\E C_k=\frac1k.}
\]
$C_k(C_k-1)$ 计数有序的两个不同 $k$ 循环。它们必须不相交。若 $2k\le n$，其期望为
\[
\binom nk\binom{n-k}k
\frac{[(k-1)!]^2(n-2k)!}{n!}=\frac1{k^2};
\]
若 $2k>n$，两个这样的循环不能共存，阶乘矩为零。因此
\[
\boxed{\Var C_k=\frac1k-\frac{\ind\{2k>n\}}{k^2}.}
\]
类似地，$j\ne k$ 时，若 $j+k\le n$，选取两个不相交子集并分别组成循环，得到 $\E(C_jC_k)=1/(jk)$；若 $j+k>n$，乘积恒为零。于是
\[
\boxed{\Cov(C_j,C_k)=-\frac{\ind\{j+k>n\}}{jk},\quad j\ne k.}
\]
\pitfall 协方差为零不等于独立。例如所有循环长度还必须满足 $\sum_k kC_k=n$。阶乘矩把“不同循环”的计数要求直接写入了随机变量。
